% Exercise 4.4
\begin{itemize}
    \item[(a)] Since $X$ is uniformly distributed on $[0, 1]$ and we are always
        using observations within 10\% of the range of $X$ closest to the test
        observation, we will use $1/10$ of the available observations to make the
        prediction, on average.
    \item[(b)] In this case, we would be using $1/10 \cdot 1/10 = 1/100$ of the
        available observations to make the prediction.
    \item[(c)] We would be using $1/{10}^{100}$ of the available observations to
        make the prediction.
    \item[(d)] In general, for a uniformly distributed random variable $X$ with
        $p$ features, where we restrict the number of observations used to a
        fraction $q$ of the given observations in the range $[0, 1]$ for each
        feature, we would be restricting our usage to $q^p$ as a fraction of
        the available observations. Thus, for $q$ relatively small, as $p$ grows
        larger, the fraction of the observations used shrinks dramatically, so
        that there are very few training observations ``near'' any given test
        observation.
    \item[(e)] If the hypercube has sides of length $l$, then it contains on
        average a fraction $l^p$ of the training observations, so that we require
        $l^p = 0.1$, i.e., $l = {0.1}^{1/p}$. Thus the length of each side of the
        hypercube is $0.1$ for $p = 1$, ${0.1}^{1/2} \approx 0.316$ for $p = 2$,
        and ${0.1}^{1/100} \approx 0.977$ for $p = 100$. This reflects the problem
        which we discussed in (d): to contain 10\% of the training observations
        when $p = 100$, the hypercube must span almost the entire range of each
        feature, so that the ``nearest'' neighbours are not really near to the
        test observation at all.
\end{itemize}
